\section{Exercise and session implementation}
Migration 0005 adds account-scoped \texttt{lab\_runs}, \texttt{task\_sessions} and
\texttt{xp\_awards}. A partial unique index permits one running session per account;
a unique account/pack/node key prevents duplicate rewards. Existing progress and
missions are retained. A native instance records its pack fingerprint, randomized
public fixture, private expected answer, attempts and verified time. A pack change
invalidates its prior instance evidence; a superseded instance cannot be submitted.

The \texttt{lab} verifier accepts \texttt{exercise} equal to logging, subnet,
access-policy, integrity, timeline or case. Case tasks additionally require a bundled
\texttt{case\_id}. Typed equality prevents JSON true from substituting for integer 1.
Account-scoped lookups, transactions and completion guards apply to the browser,
remote CLI, manual marking, force/backfill and split operations. Verified nodes
cannot be completed by splitting them into manually marked children.

The \texttt{ctf} verifier accepts a provider from the external capability catalog,
a required \texttt{challenge}, and optional \texttt{all\_of} string IDs. It checks an
HMAC-SHA256 administrator observation with an exact integer account ID, provider,
solved-ID list and timezone-aware observation time. Receipt freshness is five minutes
with five seconds of future-clock tolerance. The tracker reads a mounted receipt;
it has neither a Docker socket nor network membership in a vulnerable target.
Administrator access to the signing key is trusted. This is not tamper-proof exam
proctoring, and shared game state cannot establish which person performed a solve.

The existing authenticated \texttt{POST /api/command} accepts the session and lab CLI
arguments. No browser-supplied score, elapsed time or solved boolean is accepted.
Scheduler checkpoints use the existing account notification settings. The browser's
countdown is a preview, refreshed from the server every 30 seconds while visible.
Rendering the clock preserves its action buttons so keyboard focus is not lost each
second. Tracks contains the searchable source-derived catalog; assessment metadata
appears progressively inside challenge details.

\subsection{Regression contracts}
Exercise tests derive answers from the public generated fixture, not stored expected
values. They cover positive and negative controls, no reward for wrong answers,
account boundaries, typed comparisons, stale fingerprints, superseded instances,
manual bypass attempts, continuous timing, end-without-reward, declining rewards and checkpoint deduplication.
Receipt tests cover malformed solved lists, signatures, account mismatches, stale
and future timestamps, multi-challenge requirements and reward idempotency. These
are bounded correctness tests, not measurements of learning effectiveness.

\subsection{Reading a challenge and ending an attempt}
\path{GET /api/challenges/{reference}} resolves the caller's installed pack and node,
returns the node's level/status, its task session and the current public exercise.
\texttt{labs.current} calls the same ownership, fingerprint and expiry checks as
\texttt{labs.view}; it never calls \texttt{labs.start} or \texttt{timing.start}.
An invalid instance is reported through \texttt{exercise\_error} while the task's
instructions remain available. Reading a workspace must not create random fixtures,
resume timing or grant completion.

\texttt{timing.end} records elapsed seconds and a finished timestamp without
calling the completion or reward path. \texttt{timing.finish} awards XP only when
the prior row was running and the evidence kind is a qualifying verifier. Finished
and legacy paused attempts cannot preserve a higher reward for later submission.
The old paused value remains in migration 0005 for compatibility; no new migration
rewrites learner history. \texttt{lab start --minutes} nests \texttt{timing.start}
inside the exercise transaction, so a failed setup rolls back the new timer too.
